\section{Concluding Remarks}

This study designed, developed, and evaluated an IoT-enabled fleet-monitoring and fuel-anomaly-detection system for delivery-vehicle operations. The prototype integrates vehicle-sensor acquisition, GPS tracking, resilient multi-channel wireless transmission, cloud storage, dashboard and human-machine-interface visualisation, context-aware anomaly detection, and closed-loop alert delivery into a single operational system evaluated against the five specific objectives set at the outset.

The measured outcomes support the objectives as met. For SO1, the filtered fuel signal reached $1.04\%$ mean absolute percentage error and $3.00\%$ maximum error across the $13$--$98\%$ tank-level range against a legally calibrated pump-volume reference, with a residual-noise band below $0.1\%$ across every reported tank band; the moving-average stage of the fuel-preprocessing cascade was retained after a two-part comparison against a Kalman-filter alternative, comprising the four-axis corpus comparison that showed the two smoothers effectively tied on noise and signal preservation but decisively separated in favour of the moving average on downstream detection quality, and a supplementary live cross-scenario ablation in which the moving average again reduced consecutive-sample noise by a factor of four-to-five across city, climb, highway, idle, and refuel windows, whereas the Kalman filter was near-passthrough on the raw sender stream. For SO2, the multi-channel communication stack (LoRa~$\rightarrow$~LTE~$\rightarrow$~GSM~2G~$\rightarrow$~WiFi~$\rightarrow$~offline buffer) sustained bounded per-channel latency (LoRa and WiFi medians near seven seconds, LTE as a delayed-delivery path, GSM~2G as an automatic fringe fallback) and confirmed that approximately two-thirds of the field corpus reached the backend through the buffered or replayed path, so a real-time-only design would have lost most of the operational data. For SO3, the GPS pipeline delivered a $5.52$~s median inter-sample cadence and $90.21\%$ fix availability across $5{,}181$ field-test rows, well inside the objective's targets. For SO4, the context-aware Isolation Forest Model~C reached $92.9\%$ precision, $92.9\%$ recall, and $1.09\%$ false-positive rate under strict leave-one-trip-out evaluation, and a formal ablation with three alternative detector configurations (physics tripwire alone, Matrix Profile alone, IF~$+$~MP hybrid) confirmed that the Isolation Forest is the strongest single detector on this corpus, that Matrix Profile is best retained as an audit-only signal rather than as a co-equal union gate, and that the physics tripwire remains a defensible deterministic operational floor. For SO5, the alert engine, the embedded HMI, and the administrative dashboard together sustained a $95\%$ closed-loop alert-resolution rate on a $100$-alert field audit and an $81.9$-point System Usability Scale response from the intended administrator group.

The statistical layer sitting under those point estimates strengthens the conclusion rather than merely restating it. McNemar's paired-comparison test rejected the null of equal error rates between the Isolation Forest and every other detector configuration at $p < 0.001$, the bootstrap $95\%$ confidence intervals for the Isolation Forest and the hybrid did not overlap on $F_1$ or precision, and the Wilson score intervals for packet delivery, GPS fix availability, and alert-resolution proportion each sat above their respective objective targets with a statistical margin. Where the evidence did not support a strong statistical claim (the Wilcoxon test on the moving-average versus Kalman-filter per-trip $F_1$ differences at $p = 0.313$, and the one-sample Wilcoxon on the four SUS respondents at $p = 0.0625$), the honest reading is documented alongside the point estimate rather than glossed over, and a larger evaluation partition is recommended in the corresponding case.

The scalability analysis reported in Section~\ref{sec:so_commercial_scalability} of Chapter~\ref{ch:result_discuss} extends the study findings from the two-vehicle thesis pilot to representative commercial-fleet sizes, evaluating the prototype at $10$, $100$, $500$, and $1{,}000$ vehicles against the three closed-form load-axis ceilings derived in Section~\ref{sec:scalability_analysis}: backend ingest, LoRa gateway occupancy, and hot-table storage growth. That commercial-scale projection is the operational bridge between the measured prototype metrics and a fleet-scale rollout plan, and it identifies the LoRa gateway occupancy as the first quantitative constraint that binds at typical depot fleet sizes and the archival cadence as the storage constraint that dominates once daily hot-table growth exceeds the dashboard query-latency envelope.

Taken together, the study demonstrates that combining resilient embedded telemetry, contextual feature engineering, context-aware time-series and tree-based detectors, cloud-side processing, and multi-surface visualisation can support practical fuel monitoring and fleet supervision under real Philippine operating conditions. It also shows that fuel anomaly detection should not rely only on raw fuel-level changes, since normal driving conditions, idling, traffic, refuelling, and sensor noise can resemble abnormal behaviour; trip context, per-vehicle baselines, communication resilience, and human-readable alert reasoning are all load-bearing for a system that must be trusted by an operator. The remaining boundaries of the deployed prototype (parked-depot siphoning detection, larger real-anomaly corpus, hardware-anchored unit identification, and stronger real-time notification) are documented as the future-work directions in Sections~\ref{sec:concluding_recommendations} and~\ref{sec:concluding_future_prospects} below.

\section{Technical Limitations}

The prototype met the numerical targets of the study, but several implementation limits remain important for interpreting the results and planning deployment beyond thesis scale. Table~\ref{tab:technical_limitations} summarizes the main technical limitations, their causes, and the corresponding future-work direction.

\begin{table}[!htbp]
\centering
\caption{Technical Limitations and Future Work}
\label{tab:technical_limitations}
\scriptsize
\setlength{\tabcolsep}{3pt}
\renewcommand{\arraystretch}{1.15}
\begin{tabularx}{\textwidth}{|p{0.25\textwidth}|X|X|}
\hline
\textbf{Limitation} & \textbf{Why it happened} & \textbf{How future work fixes it} \\
\hline
No dipstick calibration & The partner fleet operator did not permit tank opening or invasive tank-volume measurement during field testing. & Conduct controlled laboratory tank calibration or partner-approved calibration using a supervised tank-opening procedure. \\
\hline
Small real anomaly count & Actual theft, leak, and siphoning events were unavailable during the thesis field window, so labels were primarily controlled injections. & Extend SGSA fleet deployment over a longer period and collect confirmed operational anomaly cases for retraining and validation. \\
\hline
LTE latency & Cellular modem attachment, HTTP actions, TLS negotiation, and network conditions can delay some telemetry uploads. & Use persistent bearer sessions, optimized modem power design, and a stronger LTE Cat-1 or equivalent modem configuration. \\
\hline
GSM 2G degradation & Carrier availability and legacy-network phaseout reduce the reliability of GSM-only fallback in some areas. & Expand LTE Cat-1, NB-IoT, LoRaWAN, or other modern fallback channels and treat 2G only as transitional support. \\
\hline
\end{tabularx}
\normalsize
\end{table}

The current deployed fuel-anomaly detector is strongest for in-trip and moving-vehicle anomalies because it uses contextual movement, speed, RPM, engine load, route state, and driver-baseline deviation to decide whether fuel loss is plausible. Parked-depot siphoning is a remaining limitation: when the vehicle is stationary and the engine is off, the production feature set has little movement or load context to compare against. Future work should add a separate stationary-state detector based on long-duration fuel-level monitoring so overnight or depot siphoning can be handled without weakening the moving-vehicle Model~C threshold.

\section{Contributions}

The interrelated \index{contributions} contributions and supplements that have been developed by the author(s) in this study are listed as follows. Only those that are unique to the authors' work are included.

\begin{itemize}
  \item the development of an integrated fleet-vehicle telemetry monitoring prototype that combines fuel-level sensing, GPS tracking, speed monitoring, trip logging, HMI alerts, and cloud dashboard visualization;

  \item the implementation of a multi-path communication approach using LoRa, LTE, GSM 2G, WiFi, and SPIFFS offline replay to support telemetry transmission under different operating and network conditions;

  \item the creation of a context-aware fuel anomaly detection pipeline using Isolation Forest, route and driving-context features, per-asset fuel-economy deviation, and comparison against baseline detectors such as static thresholding, IQR-based methods, Local Outlier Factor, Matrix Profile gating, and hybrid detection;

  \item the construction of controlled-injection validation trips and field-test telemetry records for evaluating fuel anomaly detection, trip monitoring, dashboard behavior, alert generation, and system reliability;

  \item the implementation of backend alert logic and HMI-side immediate warnings for operational events such as overspeeding and low fuel, improving response visibility for both the driver and fleet administrator;

  \item the development of dashboard analytics for fleet-level monitoring, including trip summaries, anomaly heatmaps, alert trends, fuel behavior, and overspeeding events by fleet vehicle.
\end{itemize}

\section{Recommendations}
\label{sec:concluding_recommendations}

Based on the results of the study, the researchers recommend further improvement of the system before wider field deployment. First, the fuel anomaly detection model should be evaluated using more real-world trips from different fleet vehicles, routes, load conditions, traffic levels, and driving behaviors. Although the controlled-injection and field-validation results showed promising results, additional field data would make the model more robust and reduce dependence on predefined trip patterns.

Second, the fuel detection logic should continue to distinguish between true abnormal fuel loss and normal operational events such as idling, slopes, refueling, sensor fluctuation, and traffic congestion. The results showed that some anomaly types are easier to detect than others, so future versions should improve sensitivity to short-duration events without increasing false alarms.

Third, the system should improve trip-distance and odometer handling by maintaining reliable server-side trip accumulation and reducing dependence on unavailable vehicle odometer readings. Since some vehicle OBD-II odometer data may be inaccessible, trip distance should be computed using a defensible combination of GPS, speed, and onboard integrated distance, with safeguards against corrupted or non-monotonic readings.

Fourth, the dashboard and alert system should be tested with more administrators and drivers in real fleet operations. The acceptance survey showed positive usability, but a larger group of users would provide better feedback on alert clarity, dashboard navigation, and operational usefulness.

Lastly, the researchers recommend improving deployment reliability by adding stronger logging, automated data-quality checks, and clearer handling of offline telemetry. These additions would help ensure that delayed, buffered, or partially missing data does not distort trip summaries, anomaly counts, or dashboard analytics.

Additionally, the researchers recommend hardening the unit-identification chain before broader fleet rollout. The current identification scheme is anchored at backend-issued UUIDs (\texttt{truck\_id}) rather than at hardware-rooted identities, so a swapped ESP32 controller inheriting the same \texttt{truck\_id} through NVS restore would produce telemetry indistinguishable from the original device. Three concrete extensions are recommended: (1) bind each \texttt{truck\_id} to an immutable hardware anchor such as the ESP32 factory MAC or eFuse-based secure device ID, so that a controller swap is detected on first check-in rather than silently inheriting the previous vehicle's telemetry stream; (2) tag every telemetry payload with a firmware provenance stamp (Git commit SHA plus build timestamp), so that rolling upgrades can be distinguished on the same physical unit; and (3) surface sensor-level identifiers (OBD-II VIN, GPS module UID, LoRa EUI, cellular IMEI) as part of the unit record on the operator dashboard, so that data-quality issues can be attributed to a specific sensor without a physical inspection.

\section{Future Prospects}
\label{sec:concluding_future_prospects}

There are several prospects that may be extended for further studies. The suggested topics are listed as follows:

\begin{enumerate}
  \item Integration of additional vehicle parameters not yet covered by the prototype, such as battery voltage, diagnostic trouble codes, tire-pressure data, and brake-system indicators, to improve anomaly detection and vehicle-health monitoring.

  \item Expansion of the anomaly detection model using a larger real-world fleet dataset across more routes, vehicles, fuel behaviors, and operating conditions.

  \item Development of a mobile driver application with live trip status, route guidance, alert acknowledgement, and driver-side incident reporting.

  \item Development of route-deviation and geofence-based supervision features to notify administrators when a fleet vehicle leaves an assigned corridor, enters a restricted service area, or remains stopped outside an expected location.

  \item Improvement of anomaly explainability, where each fuel anomaly alert includes the likely reason for the flag, such as sudden fuel drop, abnormal fuel-per-kilometer behavior, prolonged idling loss, or route-context mismatch.

  \item Deployment of a stronger real-time notification system using SMS, push notifications, or messaging integrations for urgent events such as fuel theft, low fuel, overspeeding, or route deviation.

  \item Evaluation of the system in long-term fleet operation to measure its effect on fuel accountability, response time, operational visibility, and maintenance planning.

  \item Extension of unit identification and hardware-provenance tracking so that fleet-scale rollouts remain auditable, by binding each vehicle's \texttt{truck\_id} to an immutable hardware anchor (ESP32 factory MAC or eFuse-based secure device ID) instead of a backend-issued UUID stored in NVS alone, tagging every telemetry payload with a firmware provenance stamp (Git commit SHA and build timestamp), and surfacing sensor-level identifiers such as OBD-II VIN, GPS module UID, LoRa EUI, and cellular IMEI on the dashboard so that data-quality issues can be attributed to a specific sensor without a physical inspection.
\end{enumerate}
